% Laget automatisk av elevark/lag_elevark.py fra sim/enheter.html. Endringer her blir overskrevet.
\documentclass[11pt]{article}
\usepackage{elevark}
\begin{document}
\elevarkhode{Fysikk 1 · Introduksjon til fysikk}{Enheter og standardform}{https://lektorbo.no/sim/enheter.html}{Simuleringen}

I fysikk regner vi med SI-enheter: meter, kilogram og sekund. Her får du tilfeldige oppgaver der du gjør om enheter, skriver tall på standardform og teller gjeldende siffer. Du får svaret og en forklaring med en gang.

\section*{Oppgaver}

\begin{oppgaveliste}
\oppgave \fase{test}{Øv}Gjør ti oppgaver av hver type. Hvilken type er vanskeligst for deg?
\svar
\oppgave \fase{forklar}{Forklar}Hvorfor deler vi på 3,6 når vi gjør om fra km/h til m/s? Vis utregningen med 1~km = 1000~m og 1~h = 3600~s.
\svar
\oppgave \fase{forklar}{Forklar}En elev mener at 1~m² = 100~cm². Hva er feil? Tegn en kvadratmeter delt opp i kvadratcentimeter.
\svar
\oppgave \fase{test}{Regn}En bil kjører 2,50~km på 3 minutter og 12 sekunder. Regn ut gjennomsnittsfarten i m/s og i km/h, med riktig antall gjeldende siffer.
\svar
\oppgave \fase{utfordring}{Utfordring}Et hår er omtrent 70~µm tykt, og et atom har en diameter på omtrent 0,1~nm. Hvor mange atomer kan ligge etter hverandre over tykkelsen på håret? Svar på standardform.
\svar
\end{oppgaveliste}

\end{document}
